\section{Conclusion}
\label{sec:conclusion}

Large sustainability-report collections require a clear path from narrative claims to filed evidence and company-level summaries. GRASP provides this path through structured BRSR storage, filing-derived labels, principle classification, language checks, numerical review, and a four-part Transparency Score. Passages and figures remain linked to the same filing. Label sources remain visible during evaluation, and every score component has an explicit calculation. The resulting workflow combines passage review, company scoring, and sector comparison in one service.

The recovered classifier agrees with weak labels on 93.00\% of test passages (95\% CI: 92.04--93.89\%) and reaches macro-F1 of 0.9158. The test set contains 6,201 passages from 175 companies excluded from supervised training. Agreement is 63.35\% for explicitly tagged narratives, 95.15\% for concept-mapped narratives, and 90.72\% for principle notes. The fixed text-vagueness detector reaches F1 of 0.7436 on 80 development examples. The Transparency Score has Pearson correlation $r=0.6314$ with 20 historical human ratings (95\% CI: 0.2683--0.8519). These findings describe principle assignment, language screening, and score comparison under the evaluation conditions in Section~\ref{sec:evaluation}.

The saved outputs cover 1,164 companies across 72 recorded sectors. Results move through four levels: one worked company, selected companies within a sector, the complete financial-sector group, and all sectors. Financial-company scores range from 29.83 to 67.56. Among the 47 sectors containing at least five companies, Wearing Apparel has the lowest mean, 42.33, and Warehousing and Transport Support has the highest, 54.55. Across the eight companies with paired-year scores, the mean rises by 5.23 points, from 45.17 to 50.40. Six scores rise and two fall. Component changes explain these movements within the available group. Company, sector, and year views make the saved results easy to inspect from both detailed and summary perspectives.

This work can be extended using multiple levels. At the data level, collection can include annexures, charts, and additional annual filings. At the modelling level, expanded principle-labelled datasets and blind human-rating studies can support new model comparisons. Adaptation using declared training-company text can be evaluated on newly collected company groups. Business descriptions can support more detailed sector assignment and indicator-applicability rules. At the numerical-review level, extensions can compare measure definitions, units, periods, and business boundaries before generating review records. At the longitudinal level, additional annual reports can support repeated comparisons of the same companies under a recorded score version. At the framework level, the filing-based approach can be studied for international sustainability standards through explicit reporting-item mappings and separate evaluation in each setting.
